%%═══════════════════════════════════════════════════════════════════
%% Lecture 22 — Metropolis Monte Carlo
%% 77315 — Computational Physics
%%═══════════════════════════════════════════════════════════════════

\section{Lecture 22 — Metropolis Monte Carlo}\label{sec:lec22}
\hrule\vspace{1.5em}

%%──────────────────────────────────────────────────────────────────
\subsection*{Overview}
The Metropolis algorithm generates a Markov chain whose
\emph{stationary distribution} is an arbitrary target
$w(\mathbf{x})$.  It avoids the need to normalise $w$ or to
draw independent samples.  We prove the key property of
\emph{detailed balance}, discuss convergence, and apply the
method to particle transport simulation (mean free path,
random direction sampling).

%%──────────────────────────────────────────────────────────────────
\subsection{Markov Chains and Detailed Balance}\label{subsec:lec22:db}

A Markov chain is a sequence of states $\mathbf{x}_0,\mathbf{x}_1,\ldots$
where each $\mathbf{x}_{k+1}$ depends only on $\mathbf{x}_k$
(the Markov property).  The chain converges to stationary
distribution $w$ if for all pairs of states $\mathbf{x},\mathbf{y}$:
\begin{equation}
  w(\mathbf{x})\,T(\mathbf{x}\to\mathbf{y})
  = w(\mathbf{y})\,T(\mathbf{y}\to\mathbf{x}),
  \label{eq:lec22:detailed_balance}
\end{equation}
where $T(\mathbf{x}\to\mathbf{y})$ is the transition probability.
This \emph{detailed balance} (microscopic reversibility) implies
$w$ is preserved by the chain.

%%──────────────────────────────────────────────────────────────────
\subsection{Metropolis Algorithm}\label{subsec:lec22:metro}

Choose a \emph{proposal distribution} $q(\mathbf{y}|\mathbf{x})$
(e.g.\ uniform in a ball of radius $\delta$ around $\mathbf{x}$).

\begin{enumerate}
  \item Start at state $\mathbf{x}$.
  \item Propose $\mathbf{y}$ from $q(\cdot|\mathbf{x})$.
  \item Compute acceptance ratio:
        \begin{equation}
          r = \frac{w(\mathbf{y})\,q(\mathbf{x}|\mathbf{y})}
                   {w(\mathbf{x})\,q(\mathbf{y}|\mathbf{x})}.
          \label{eq:lec22:acceptance}
        \end{equation}
        For symmetric proposals $q(\mathbf{x}|\mathbf{y})=q(\mathbf{y}|\mathbf{x})$
        (Metropolis–Hastings reduces to original Metropolis):
        $r = w(\mathbf{y})/w(\mathbf{x})$.
  \item Draw $u\sim\text{Uniform}(0,1)$.
        Accept: $\mathbf{x}\leftarrow\mathbf{y}$ if $u < \min(1,r)$;
        otherwise $\mathbf{x}\leftarrow\mathbf{x}$ (stay).
  \item Discard a burn-in period (equilibration).
  \item Use remaining $\mathbf{x}_k$ as samples from $w$.
\end{enumerate}

\subsubsection*{Proof of detailed balance}
For a transition from $\mathbf{x}$ to $\mathbf{y}$ with $r<1$:
$T(\mathbf{x}\to\mathbf{y}) = q(\mathbf{y}|\mathbf{x})\cdot r
= q(\mathbf{y}|\mathbf{x})\,w(\mathbf{y})/w(\mathbf{x})$.
Then $w(\mathbf{x})\,T(\mathbf{x}\to\mathbf{y})
= q(\mathbf{y}|\mathbf{x})\,w(\mathbf{y})
= w(\mathbf{y})\,T(\mathbf{y}\to\mathbf{x})$ (since $r^{-1}\ge 1$
means the reverse transition is always accepted). $\square$

%%──────────────────────────────────────────────────────────────────
\subsection{Practical Considerations}\label{subsec:lec22:practice}

\subsubsection*{Step size $\delta$}
\begin{itemize}
  \item Too large: almost all proposals rejected; chain explores slowly.
  \item Too small: proposals always accepted but correlation length
        is very long — also slow decorrelation.
  \item Optimal acceptance rate $\approx 23$\%–$50$\% (dimension-dependent).
\end{itemize}

\subsubsection*{Autocorrelation time}
Consecutive samples $\mathbf{x}_k$, $\mathbf{x}_{k+1}$ are
correlated.  The \emph{autocorrelation time} $\tau$ measures the
number of steps between effectively independent samples.
The effective sample size is $N_{\rm eff} = N / (2\tau)$.

\nt{$w(\mathbf{x})$ need not be normalised — only the ratio $r = w(\mathbf{y})/w(\mathbf{x})$ is needed.  This is the primary advantage over direct sampling methods.}

%%──────────────────────────────────────────────────────────────────
\subsection{Particle Transport Monte Carlo}\label{subsec:lec22:transport}

\subsubsection*{Mean free path}
A particle moving through a medium of number density $n_0$
with cross-section $\sigma$ has mean free path
\begin{equation}
  \bar\ell = \frac{1}{n_0\sigma}.
  \label{eq:lec22:mfp}
\end{equation}
The path length to the next collision is exponentially distributed:
\begin{equation}
  P(\ell) = \frac{1}{\bar\ell}\,e^{-\ell/\bar\ell}.
  \label{eq:lec22:path}
\end{equation}
Sample from \eqref{eq:lec22:path} using the inverse-CDF:
$\ell = -\bar\ell\,\ln(1-u)$, $u\sim\text{Uniform}(0,1)$.

\subsubsection*{Random direction}
Sample an isotropic direction in 3D by drawing
$\cos\theta \sim \text{Uniform}(-1,1)$ and
$\phi \sim \text{Uniform}(0,2\pi)$:
\begin{equation}
  \hat{n} = (\sin\theta\cos\phi,\;\sin\theta\sin\phi,\;\cos\theta).
  \label{eq:lec22:direction}
\end{equation}

\subsubsection*{Simulation loop}
\begin{enumerate}
  \item Initialise particle position and direction.
  \item Sample path length $\ell$ from \eqref{eq:lec22:path}.
  \item Advance particle; record tallies (energy deposition, flux).
  \item Sample new direction from scattering kernel.
  \item Repeat until the particle escapes or is absorbed.
\end{enumerate}
Repeat for $N$ particles; physical observables are ensemble averages
with error $\sim N^{-1/2}$.
